\section{问题定义与可行性分析}

\subsection{项目背景}

大模型 API 的计费日趋复杂：厂商既提供附带调用额度的包月订阅，也提供按 Token 计费的 API 渠道。按量计费又区分输入、输出、缓存读取、缓存写入四种单价，并随服务档位（Priority、Flex）和上下文长度变化。个人单独购买多个订阅成本高、利用率低，因此常见做法是多人合买若干上游账号，通过 API 网关（本项目对接开源网关 CLIProxyAPI，简称 CPA）向每人分发独立密钥，共享使用。

共享带来了核算问题。月度总成本由固定订阅费和按量费用组成，成员用量相差可达数十倍，订阅周期与自然月不对齐，价格经常调整，密钥还会被重置、吊销或转手。仅靠网关统计和电子表格手工结算，会出现以下问题：
\begin{enumerate}
  \item 调用记录分散在网关的 SQLite 日志库中，缺少用户维度，无法回答“某人本月在哪些模型上花了多少钱”；
  \item 价格变化后无法说明某一期账单是按哪一版价格算出来的，账单不可复核；
  \item 密钥转手后，历史调用被错误地计入新主人名下；
  \item 分摊规则（例如用量越大单价越低的累进折扣）改动后，已经结清的月份会被“追溯”改变。
\end{enumerate}

\subsection{问题定义}

本课题要建立一个以关系数据库为核心的“大模型 API 用量统计与费用分摊系统”，其目标是：
\begin{itemize}
  \item \textbf{统一管理数据}：用户、API 密钥及其历史归属、调用记录、版本化的模型价格、上游账号成本、分摊规则、账期与账单；
  \item \textbf{调用级计价}：对每一次调用按当时生效的价格版本计算等效费用，并保存计算明细；
  \item \textbf{多维统计}：支持按用户、模型、时间、上游账号等维度统计请求数、Token 与费用；
  \item \textbf{共享成本分摊}：把每个账期的上游总成本按梯度规则公平地分到每位成员，精确到分且总额守恒；
  \item \textbf{账单可查询、可核对}：每期账单记录所用价格版本、规则快照与成本快照，关账后数据由数据库保证不可修改。
\end{itemize}

系统边界：网关 CPA 负责转发调用与管理密钥，用量库 CPAMP（网关的用量记录与价格同步组件）负责记录原始调用日志和抓取公开价格。本系统以只读方式增量读取用量库，通过网关管理接口增删密钥和查询上游额度，自身不转发任何模型调用。

\subsection{可行性分析}

\paragraph{技术可行性}
按编程助手单人每天数百次调用估算，30 人左右的共享组每月约 10 万至 50 万次调用，三年累计不超过数千万行，单机 PostgreSQL 足以承载；第 7 章的测试表明百万级数据上常用查询在 5 ms 内完成。计价与分摊均为整数运算，没有浮点误差。所用框架均为成熟的开源项目。

\paragraph{经济可行性}
系统全部基于开源软件，可与网关部署在同一台 2 核 4 GB 云服务器上，不增加硬件费用，而自动分摊可以省去每月数小时的手工对账。

\paragraph{操作可行性}
成员通过已在使用的 Telegram 机器人注册和管理密钥，Web 端用自己的密钥登录；管理员的日常操作只有“创建账期—预览—关账”三步，其余由后台任务完成。

\paragraph{法律与安全可行性}
系统不存储提示词与回复内容；密钥只保存不可逆摘要和脱敏值，上游账号凭据始终由网关保管。

\subsection{开发环境与工具}

\begin{table}[H]
\centering\small
\caption{开发与运行环境}
\begin{tabularx}{\textwidth}{@{}lX@{}}
\toprule
类别 & 选型 \\ \midrule
数据库管理系统 & PostgreSQL 17（生产可选），SQLite 3.40+（单机默认，兼容） \\
后端 & Python 3.12，FastAPI，SQLAlchemy 2.1（ORM 与 Core），Alembic 1.20（模式迁移），psycopg 3 \\
前端 & Vue 3.5，Vuetify 3，ECharts 6，Vite 6，Vitest \\
交互入口 & Web 控制台（B/S），Telegram Bot，命令行工具 \\
设计工具 & Graphviz（E-R 图、数据流图），本报告使用 XeLaTeX 排版 \\
部署 & Docker Compose，GitHub Actions 持续集成（SQLite 与 PostgreSQL 双库测试） \\
\bottomrule
\end{tabularx}
\end{table}
